\section{서론}

소형 멀티로터 UAV는 자세 추정의 주요 감지 수단으로 관성 측정 장치에 의존한다.
MEMS IMU는 3축 자이로스코프와 가속도계, 그리고 9-DOF 구성의 경우 지자기 센서를 결합하여
월드 프레임에 대한 롤, 피치, 요를 추정한다.
IMU 기반 추정의 근본적인 문제는 자이로스코프가 적분을 통해 드리프트를 누적하는 반면,
가속도계는 진동 및 선형 가속도로 오염된 절대 기울기 기준값을 제공한다는 점이다~\cite{woodman2007}.
이러한 신호들의 상호 보완적인 주파수 특성을 활용하는 센서 융합 알고리즘은
단순한 상보 필터부터 완전한 확률적 추정기에 이르기까지 광범위하게 연구되어 왔다.

비행 컨트롤러 측면에서 PX4~\cite{px4docs}와 ArduPilot~\cite{ardupilotdocs} 같은
오픈소스 자동 조종 프레임워크는 상태 추정을 위한 온보드 EKF와 결합된
계단식 비례-미분 자세 및 각속도 제어기(PPID)를 구현한다.
추정 알고리즘의 선택이 폐루프 자세 제어 성능에 미치는 영향을 이해하는 것은
임베디드 배포를 위한 하드웨어 및 소프트웨어 스택을 선택하는 UAV 설계자에게 실질적으로 중요하다.

본 논문은 두 가지 독립적이면서도 연관된 연구를 통해 문제의 두 수준을 모두 다룬다.
첫 번째 연구는 노이즈 모델링, 지터 분석, MuJoCo 시뮬레이션 기준값을 활용한
여섯 가지 필터 비교 평가를 통해 MPU-9250 IMU를 특성화한다.
두 번째 연구는 SK450 쿼드로터의 하드웨어 구조를 분석하고
PX4/ArduPilot의 PPID 및 EKF 제어 파이프라인을 MuJoCo 시뮬레이션으로 재현한다.
논의 섹션에서는 IMU 필터 비교가 비행 컨트롤러 추정기 설계에 주는 시사점을 교차 분석한다.

본 논문의 기여는 다음과 같다:
\begin{itemize}
    \item 아두이노급 임베디드 하드웨어에서 MPU-9250의 노이즈 특성화 및 지터 분석.
    \item MuJoCo 시뮬레이션 기준값과 실시간 하드웨어 평가를 사용한 여섯 가지 자세 추정 필터의 정량적 비교.
    \item SK450 + CUAV 7 Nano 하드웨어 구조 및 제어 파이프라인 분석.
    \item PX4 및 ArduPilot에 문서화된 PPID 계단식 제어기 및 EKF의 MuJoCo 시뮬레이션 구현 및 검증.
    \item IMU 필터 선택 기준을 실용적인 비행 컨트롤러 설계 시사점으로 연결하는 교차 논의.
\end{itemize}

본 논문의 나머지 구성은 다음과 같다.
2절에서 관련 연구를 검토하고,
3절에서 IMU 상태 추정 연구를 제시하며,
4절에서 UAV 플랫폼 및 시뮬레이션을 설명하고,
5절에서 교차 연구 시사점을 논의하며,
6절에서 결론을 맺는다.
